\documentclass[11pt]{article}
\usepackage[sans]{sim-oficial}
\input{datos-proponente}
\newcolumntype{R}{>{\raggedleft\arraybackslash}p{2.1cm}}
\renewcommand{\arraystretch}{1.18}
% Statement heading: #1 title, #2 period line
\newcommand{\eeffhead}[2]{%
  \begin{center}
    {\large\bfseries LAFROX SpA}\\[3pt]
    {\bfseries #1}\\[2pt]
    #2\\[2pt]
    {\small (Cifras en millones de pesos chilenos --- MM\$)}
  \end{center}\vspace{2pt}}
\newcommand{\seccion}[1]{\multicolumn{4}{@{}l}{\textbf{#1}}\\}
\newcommand{\fil}[4]{\hspace{1em}#1 & #2 & #3 & #4\\}
\newcommand{\tot}[4]{\textbf{#1} & \textbf{#2} & \textbf{#3} & \textbf{#4}\\}
\newcommand{\cabeera}{\toprule & \textbf{2025} & \textbf{2024} & \textbf{2023}\\\midrule}
\newcommand{\nota}[2]{\par\medskip\noindent\textbf{#1}\ #2}
\begin{document}
\pagestyle{fancy}

%--------------------------------------------------------------- portada
\thispagestyle{fancy}
\vspace*{3.5cm}
\begin{center}
  {\Huge\bfseries LAFROX SpA}\\[26pt]
  {\Large\bfseries ESTADOS FINANCIEROS}\\[8pt]
  {\large Al 31 de diciembre de 2025, 2024 y 2023}\\[8pt]
  {\large con Informe de los Auditores Independientes}\\[30pt]
  {\small (Cifras en millones de pesos chilenos --- MM\$)}
\end{center}
\vfill
\noindent\textbf{CONTENIDO}\par\smallskip
\noindent\begin{tabularx}{\linewidth}{@{}X r@{}}
  Informe de los Auditores Independientes & 2\\
  Estados de Situación Financiera & 3--4\\
  Estados de Resultados Integrales por Función & 5\\
  Estados de Flujos de Efectivo & 6\\
  Estados de Cambios en el Patrimonio & 7\\
  Notas a los Estados Financieros & 8\\
  Indicadores financieros & 9\\
\end{tabularx}
\par\medskip\noindent{\footnotesize NIIF: Normas Internacionales de Información Financiera. Estados financieros simulados con fines académicos (Art.~34 y Art.~39\textdegree, numeral~2, de las Bases Administrativas). Cifras y auditor ficticios.}
\clearpage

%--------------------------------------------------------------- informe auditor
\begin{center}{\large\bfseries INFORME DE LOS AUDITORES INDEPENDIENTES}\end{center}
\noindent Señores\\Accionistas y Directores\\\LFXempresa

\smallskip
\noindent\textbf{Opinión}

\noindent Hemos auditado los estados financieros de \LFXempresa{} (la ``Sociedad''), que comprenden los estados de situación financiera al 31 de diciembre de 2025, 2024 y 2023 y los correspondientes estados de resultados integrales, de cambios en el patrimonio y de flujos de efectivo por los años terminados en esas fechas, y las notas a los estados financieros.

En nuestra opinión, los estados financieros adjuntos presentan razonablemente, en todos sus aspectos significativos, la situación financiera de \LFXempresa{} al 31 de diciembre de 2025, 2024 y 2023, y los resultados de sus operaciones y sus flujos de efectivo por los años terminados en esas fechas, de acuerdo con las Normas Internacionales de Información Financiera (NIIF). \textbf{Nuestra opinión es sin salvedades.}

\noindent\textbf{Base para la opinión}

\noindent Efectuamos nuestras auditorías de acuerdo con las normas de auditoría generalmente aceptadas en Chile. Nuestras responsabilidades de acuerdo con tales normas se describen más adelante en la sección ``Responsabilidades del auditor''. Somos independientes de la Sociedad y hemos cumplido las demás responsabilidades éticas pertinentes. Consideramos que la evidencia de auditoría obtenida es suficiente y apropiada para proporcionar una base para nuestra opinión.

\noindent\textbf{Responsabilidad de la administración por los estados financieros}

\noindent La administración es responsable de la preparación y presentación razonable de estos estados financieros de acuerdo con las NIIF; esto incluye el diseño, implementación y mantención del control interno pertinente para la preparación y presentación razonable de estados financieros que estén exentos de representaciones incorrectas significativas, ya sea debido a fraude o error, y de evaluar la capacidad de la Sociedad para continuar como empresa en marcha.

\noindent\textbf{Responsabilidades del auditor}

\noindent Nuestro objetivo es obtener una seguridad razonable de que los estados financieros en su conjunto están exentos de representaciones incorrectas significativas, ya sea debido a fraude o error, y emitir un informe que contenga nuestra opinión. Una auditoría incluye examinar, sobre la base de pruebas, la evidencia que respalda los montos e informaciones revelados en los estados financieros, evaluar las políticas contables utilizadas y la razonabilidad de las estimaciones significativas efectuadas por la administración, así como la presentación general de los estados financieros.

\vspace{10pt}
\noindent Santiago, 27 de febrero de 2026

\vspace{28pt}
\noindent\rule{6.4cm}{0.5pt}\\
María Muestra Ejemplo (auditora ficticia)\\
Socia\\
\SIMauditora\\[4pt]
{\small Firma, timbre y registro de auditor simulados.}
\clearpage

%--------------------------------------------------------------- ESF activos
\eeffhead{ESTADOS DE SITUACIÓN FINANCIERA}{Al 31 de diciembre de 2025, 2024 y 2023}
\noindent\begin{tabularx}{\linewidth}{@{}X R R R@{}}
  \cabeera
  \seccion{ACTIVOS}
  \seccion{Activos corrientes}
  \fil{Efectivo y equivalentes al efectivo}{820}{640}{520}
  \fil{Deudores comerciales y otras cuentas por cobrar}{1.900}{1.640}{1.400}
  \fil{Otros activos corrientes}{430}{360}{330}
  \tot{Total activos corrientes}{3.150}{2.640}{2.250}
  \midrule
  \seccion{Activos no corrientes}
  \fil{Propiedades, planta y equipos}{1.250}{1.150}{1.060}
  \fil{Activos intangibles (desarrollo de software)}{350}{330}{290}
  \fil{Otros activos no corrientes}{200}{170}{150}
  \tot{Total activos no corrientes}{1.800}{1.650}{1.500}
  \midrule
  \tot{TOTAL ACTIVOS}{4.950}{4.290}{3.750}
  \bottomrule
\end{tabularx}

\vspace{10pt}
\noindent{\small Las notas adjuntas forman parte integral de estos estados financieros.}
\clearpage

%--------------------------------------------------------------- ESF pasivos
\eeffhead{ESTADOS DE SITUACIÓN FINANCIERA}{Al 31 de diciembre de 2025, 2024 y 2023 (continuación)}
\noindent\begin{tabularx}{\linewidth}{@{}X R R R@{}}
  \cabeera
  \seccion{PASIVOS Y PATRIMONIO}
  \seccion{Pasivos corrientes}
  \fil{Proveedores y otras cuentas por pagar}{780}{690}{640}
  \fil{Obligaciones laborales y previsionales}{640}{560}{500}
  \fil{Obligaciones financieras corrientes}{330}{300}{260}
  \tot{Total pasivos corrientes}{1.750}{1.550}{1.400}
  \midrule
  \seccion{Pasivos no corrientes}
  \fil{Obligaciones financieras no corrientes}{900}{750}{680}
  \fil{Otros pasivos no corrientes}{100}{90}{70}
  \tot{Total pasivos no corrientes}{1.000}{840}{750}
  \midrule
  \seccion{Patrimonio}
  \fil{Capital emitido}{1.000}{1.000}{1.000}
  \fil{Resultados acumulados}{755}{480}{242}
  \fil{Resultado del ejercicio}{445}{420}{358}
  \tot{Total patrimonio}{2.200}{1.900}{1.600}
  \midrule
  \tot{TOTAL PASIVOS Y PATRIMONIO}{4.950}{4.290}{3.750}
  \bottomrule
\end{tabularx}

\vspace{10pt}
\noindent{\small Las notas adjuntas forman parte integral de estos estados financieros.}
\clearpage

%--------------------------------------------------------------- ERI
\eeffhead{ESTADOS DE RESULTADOS INTEGRALES POR FUNCIÓN}{Por los años terminados al 31 de diciembre de 2025, 2024 y 2023}
\noindent\begin{tabularx}{\linewidth}{@{}X R R R@{}}
  \cabeera
  Ingresos de actividades ordinarias & 6.800 & 5.900 & 5.200\\
  Costo de ventas & (4.700) & (4.130) & (3.640)\\
  \tot{Margen bruto}{2.100}{1.770}{1.560}
  Gastos de administración y ventas & (1.400) & (1.115) & (1.000)\\
  \tot{Resultado operacional}{700}{655}{560}
  Resultado financiero neto & (90) & (80) & (70)\\
  Resultado antes de impuestos & 610 & 575 & 490\\
  Impuesto a las ganancias (27\,\%) & (165) & (155) & (132)\\
  \midrule
  \tot{UTILIDAD DEL EJERCICIO}{445}{420}{358}
  \bottomrule
\end{tabularx}

\vspace{10pt}
\noindent{\small Las notas adjuntas forman parte integral de estos estados financieros.}
\clearpage

%--------------------------------------------------------------- EFE
\eeffhead{ESTADOS DE FLUJOS DE EFECTIVO}{Por los años terminados al 31 de diciembre de 2025, 2024 y 2023}
\noindent\begin{tabularx}{\linewidth}{@{}X R R R@{}}
  \cabeera
  Flujo de actividades de la operación & 720 & 610 & 540\\
  Flujo de actividades de inversión & (380) & (340) & (310)\\
  Flujo de actividades de financiación & (160) & (150) & (140)\\
  \midrule
  \tot{Incremento neto de efectivo y equivalentes}{180}{120}{90}
  Efectivo y equivalentes al inicio del año & 640 & 520 & 430\\
  \midrule
  \tot{EFECTIVO Y EQUIVALENTES AL FINAL DEL AÑO}{820}{640}{520}
  \bottomrule
\end{tabularx}

\vspace{10pt}
\noindent{\small Las notas adjuntas forman parte integral de estos estados financieros.}
\clearpage

%--------------------------------------------------------------- ECP
\eeffhead{ESTADOS DE CAMBIOS EN EL PATRIMONIO}{Por los años terminados al 31 de diciembre de 2025 y 2024}
\noindent\begin{tabularx}{\linewidth}{@{}X R R R R@{}}
  \toprule
  & \textbf{Capital emitido} & \textbf{Result. acumulados} & \textbf{Result. del ejercicio} & \textbf{Total patrimonio}\\\midrule
  \multicolumn{5}{@{}l}{\textbf{Año 2025}}\\
  \hspace{1em}Saldo inicial al 01-01-2025 & 1.000 & 480 & 420 & 1.900\\
  \hspace{1em}Traspaso del resultado anterior & --- & 420 & (420) & ---\\
  \hspace{1em}Dividendos pagados & --- & (145) & --- & (145)\\
  \hspace{1em}Resultado del ejercicio & --- & --- & 445 & 445\\
  \textbf{Saldo final al 31-12-2025} & \textbf{1.000} & \textbf{755} & \textbf{445} & \textbf{2.200}\\\midrule
  \multicolumn{5}{@{}l}{\textbf{Año 2024}}\\
  \hspace{1em}Saldo inicial al 01-01-2024 & 1.000 & 242 & 358 & 1.600\\
  \hspace{1em}Traspaso del resultado anterior & --- & 358 & (358) & ---\\
  \hspace{1em}Dividendos pagados & --- & (120) & --- & (120)\\
  \hspace{1em}Resultado del ejercicio & --- & --- & 420 & 420\\
  \textbf{Saldo final al 31-12-2024} & \textbf{1.000} & \textbf{480} & \textbf{420} & \textbf{1.900}\\
  \bottomrule
\end{tabularx}

\vspace{10pt}
\noindent{\small Las notas adjuntas forman parte integral de estos estados financieros.}
\clearpage

%--------------------------------------------------------------- Notas
\begin{center}{\large\bfseries NOTAS A LOS ESTADOS FINANCIEROS}\\[2pt]Al 31 de diciembre de 2025, 2024 y 2023 (resumen)\end{center}
\nota{Nota 1. Entidad.}{\LFXempresa{} es una sociedad por acciones constituida en 2012, con casa matriz en Valparaíso y sede central en Santiago. Su actividad es el desarrollo, integración y operación de sistemas de software y servicios de tecnologías de la información.}
\nota{Nota 2. Bases de preparación.}{Los estados financieros se preparan de acuerdo con las NIIF, bajo el criterio de devengado y en pesos chilenos, redondeados a millones.}
\nota{Nota 3. Reconocimiento de ingresos.}{Los ingresos por desarrollo se reconocen según el avance de los hitos contractuales; los de operación y soporte, mensualmente conforme al servicio prestado.}
\nota{Nota 4. Propiedades, planta y equipos e intangibles.}{Se valoran al costo menos depreciación y amortización acumuladas. El desarrollo de software capitalizable se amortiza en cuatro años.}
\nota{Nota 5. Obligaciones financieras.}{Corresponden a créditos bancarios en pesos y UF, sin cláusulas de aceleración incumplidas al cierre.}
\nota{Nota 6. Dividendos.}{En 2024 y 2025 se pagaron dividendos por 120 y 145 millones, respectivamente, con cargo a las utilidades de los ejercicios anteriores.}
\nota{Nota 7. Hechos posteriores.}{No existen hechos posteriores al cierre que afecten significativamente los estados financieros.}
\clearpage

%--------------------------------------------------------------- Indicadores
\begin{center}{\large\bfseries INDICADORES FINANCIEROS}\\[2pt]Derivados de los estados financieros\end{center}
\noindent\begin{tabularx}{\linewidth}{@{}X R R R@{}}
  \toprule \textbf{Indicador} & \textbf{2025} & \textbf{2024} & \textbf{2023}\\\midrule
  Liquidez corriente (activo corriente / pasivo corriente) & 1,80 & 1,70 & 1,61\\
  Endeudamiento (pasivo total / patrimonio) & 1,25 & 1,26 & 1,34\\
  Patrimonio (MM\$) & 2.200 & 1.900 & 1.600\\
  Ventas anuales (MM\$) & 6.800 & 5.900 & 5.200\\
  Margen operacional (\%) & 10,3 & 11,1 & 10,8\\
  Rentabilidad sobre patrimonio (\%) & 20,2 & 22,1 & 22,4\\
  \bottomrule
\end{tabularx}

\medskip
\noindent El patrimonio es positivo en los tres ejercicios, conforme al requisito de solvencia financiera del Art.~34.1 de las Bases Administrativas.

\vspace{30pt}
\noindent\SIMfecha\par\vspace{26pt}
\noindent\rule{6cm}{0.5pt}\\
\LFXrepresentante, Gerente General\\\LFXempresa
\end{document}
